% (c) 2002 Matthew Boedicker <mboedick@mboedick.org> (original author) http://mboedick.org
% (c) 2003-2007 David J. Grant <davidgrant-at-gmail.com> http://www.davidgrant.ca
% (c) 2008-2015 Nathaniel Johnston <nathaniel@njohnston.ca> http://www.njohnston.ca
%
% Depending on your TeX distribution, you may need to download the revnum and longtable packages for this template to work!
%
%This work is licensed under the Creative Commons Attribution-Noncommercial-Share Alike 2.5 License. To view a copy of this license, visit http://creativecommons.org/licenses/by-nc-sa/2.5/ or send a letter to Creative Commons, 543 Howard Street, 5th Floor, San Francisco, California, 94105, USA.

\documentclass[letterpaper,11pt]{article}
\newlength{\outerbordwidth}
\pagestyle{empty}
\raggedbottom
\raggedright
\usepackage{array}
\usepackage[svgnames]{xcolor}
\usepackage{enumerate}
\usepackage{framed}
\usepackage{longtable}
\usepackage{revnum}
\usepackage[colorlinks=true,urlcolor=blue]{hyperref}
\usepackage{tocloft}
\usepackage{etoolbox}
\robustify\cftdotfill

%-----------------------------------------------------------
%Edit these values as you see fit

\setlength{\outerbordwidth}{3pt}  % Width of border outside of title bars
\definecolor{shadecolor}{gray}{0.75}  % Outer background color of title bars (0 = black, 1 = white)
\definecolor{shadecolorB}{gray}{0.93}  % Inner background color of title bars


%-----------------------------------------------------------
%Margin setup

\setlength{\evensidemargin}{-0.25in}
\setlength{\headheight}{0in}
\setlength{\headsep}{0in}
\setlength{\oddsidemargin}{-0.25in}
\setlength{\paperheight}{11in}
\setlength{\paperwidth}{8.5in}
\setlength{\tabcolsep}{0in}
\setlength{\textheight}{9.5in}
\setlength{\textwidth}{7in}
\setlength{\topmargin}{-0.3in}
\setlength{\topskip}{0in}
\setlength{\voffset}{0.1in}
\setlength\LTleft{0.2in} % needed to make longtable full-width
\setlength\LTright{0.2in}

%-----------------------------------------------------------
%Custom commands
\newcommand{\resitem}[1]{\item #1 \vspace{-2pt}}
\newcommand{\resheading}[1]{\vspace{8pt}
  \parbox{\textwidth}{\setlength{\FrameSep}{\fboxsep}
    \begin{shaded}
\setlength{\fboxsep}{0pt}\framebox[\textwidth][l]{\setlength{\fboxsep}{4pt}\fcolorbox{shadecolorB}{shadecolorB}{\textbf{\sffamily{\mbox{~}\makebox[6.762in][l]{\large #1} \vphantom{p\^{E}}}}}}
    \end{shaded}
  }\vspace{-5pt}
}

% the next four commands allow for the \ressubheading environment to be 1, 2, 3, or 4 subrows, depending on which command you use. This is admittedly hack-ish, and should probably be replaced by a single more flexible command (with optional arguments) in the future
\newcommand{\ressubheading}[4]{
\begin{tabular*}{6.5in}[t]{l@{\cftdotfill{\cftsecdotsep}\extracolsep{\fill}}r}
		\textbf{#1} & #2 \\
		\textit{#3} & \textit{#4} \\
\end{tabular*}\vspace{-6pt}}
\newcommand{\ressubheadingb}[6]{
\begin{tabular*}{6.5in}[t]{l@{\cftdotfill{\cftsecdotsep}\extracolsep{\fill}}r}
		\textbf{#1} & #2 \\
		\textit{#3} & \textit{#4} \\
		\textit{#5} & \textit{#6} \\
\end{tabular*}\vspace{-6pt}}
\newcommand{\ressubheadingc}[8]{
\begin{tabular*}{6.5in}[t]{l@{\cftdotfill{\cftsecdotsep}\extracolsep{\fill}}r}
		\textbf{#1} & #2 \\
		\textit{#3} & \textit{#4} \\
		\textit{#5} & \textit{#6} \\
		\textit{#7} & \textit{#8} \\
\end{tabular*}\vspace{-6pt}}
\newcommand\foo[9]{%
    \def\tempb{#2}%
    \def\tempc{#3}%
    \def\tempd{#4}%
    \def\tempe{#5}%
    \def\tempf{#6}%
    \def\tempg{#7}%
    \def\temph{#8}%
    \def\tempi{#9}%
    \foocontinued
}
\newcommand\foocontinued[7]{%
    % Do whatever you want with your 9+7 arguments here.
}

\newcommand{\ressubheadingd}[1]{
	\def\argten{#1}%
	\ressubheadingdb
}
\newcommand{\ressubheadingdb}[9]{
\begin{tabular*}{6.5in}[t]{l@{\cftdotfill{\cftsecdotsep}\extracolsep{\fill}}r}
		\textbf{\argten} & #1 \\
		\textit{#2} & \textit{#3} \\
		\textit{#4} & \textit{#5} \\
		\textit{#6} & \textit{#7} \\
		\textit{#8} & \textit{#9} \\
\end{tabular*}\vspace{-6pt}}
%-----------------------------------------------------------


\begin{document}

{\large \begin{tabular*}{7in}{l@{\extracolsep{\fill}}r}
\textbf{\LARGE Luis M. B. Varona}\smallskip & \textbf{\today}\smallskip \\
Mount Allison University & luismbvarona@protonmail.com \\
62 York St., Sackville, New Brunswick, Canada \ E4L 1E2 & \hyperref{https://luismbvarona.com}{}{}{luismbvarona.com} \\
\end{tabular*}}
\\


%%%%%%%%%%%%%%%%%%%%%%%%%%%%%%
\resheading{Education}
%%%%%%%%%%%%%%%%%%%%%%%%%%%%%%
\begin{itemize}
\item
	\ressubheading{Mount Allison University}{Sackville, NB}{Bachelor of Arts with First Class Honours}{2022 -- 2027}
	\begin{itemize}
		\resitem{Triple Major: Mathematics (Hons.), Political Science, and Computer Science \& Economics}
        \resitem{Double Minor: Canadian Public Policy and International Politics}
        \resitem{Thesis: Computational Methods for Modelling Continuous-Time Quantum Walks (May 2026)}
	\end{itemize}
\end{itemize}


%%%%%%%%%%%%%%%%%%%%%%%%%%%%%%
\resheading{Research Interests}
%%%%%%%%%%%%%%%%%%%%%%%%%%%%%%
\par Combinatorial graph algorithms, graph structure of neural networks, and/or computational social science on graphs. Recent summer research centred around investigating recurrent neural networks subject to graph-theoretic constraints, and my thesis applied optimization algorithms to probe spectral properties of graphs. Previous summers of research similarly encompassed both spectral and algorithmic graph theory.\vspace{10pt}


%%%%%%%%%%%%%%%%%%%%%%%%%%%%%%
\resheading{Publications}
%%%%%%%%%%%%%%%%%%%%%%%%%%%%%%

\vspace{-0.15in}\subsection*{Peer-Reviewed Journal Articles}

\begin{revnumerate}
\item
    L.~M.~B. Varona. MatrixBandwidth.jl: Fast algorithms for matrix bandwidth minimization and recognition. \emph{Journal of Open Source Software}, 2025. doi:\hyperref{https://doi.org/10.21105/joss.09136}{}{}{10.21105/joss.09136}
\end{revnumerate}


\subsection*{Unpublished Papers}
\begin{revnumerate}
\item
	L.~M.~B.~Varona. \emph{A polynomial-time algorithm for recognizing high-bandwidth graphs}. E-print: \hyperref{https://arxiv.org/abs/2602.01755}{}{}{arXiv:2602.01755} [cs.DS], 2026. Under review at \emph{Discrete Applied Mathematics}.
    
\item
	F.~A.~Steinke and L.~M.~B.~Varona. \emph{Efficient spectral bounds on the chromatic number of Hamming, Johnson, and Kneser graph powers}. E-print: \hyperref{https://arxiv.org/abs/2601.01962}{}{}{arXiv:2601.01962} [math.CO], 2026. Under review at the \emph{Rose-Hulman Undergraduate Mathematics Journal}.

\item
	N.~Johnston, S.~Plosker, C.~Torrance, and L.~M.~B.~Varona. \emph{Generalizing the Cauchy--Schwarz inequality: Hadamard powers and tensor products}. E-print: \hyperref{https://arxiv.org/abs/2507.10327}{}{}{arXiv:2507.10327} [math.FA], 2025. Under review at \emph{Linear and Multilinear Algebra}.
\end{revnumerate}


\subsection*{Drafts in Preparation}
\begin{revnumerate}
\item
	C.~Brett and L.~M.~B.~Varona. \emph{The effects of changes in effective prices of services on municipal tax rates: evidence from New Brunswick} (working title).

\item
	W.~A.~Hunt and L.~M.~B.~Varona. \emph{Foundational models of sociopolitical change lag behind the ongoing reality of an algorithmic revolution} (working title).

\item
	L.~M.~B.~Varona. \emph{Formal language learning in bandwidth-constrained recurrent neural networks} (working title).

\item
	N.~Johnston, S.~Plosker, and L.~M.~B.~Varona. \emph{Laplacian integral and $\{-1,0,1\}$-diagonalizable graphs} (working title).
\end{revnumerate}


%%%%%%%%%%%%%%%%%%%%%%%%%%%%%%
\resheading{Presentations}
%%%%%%%%%%%%%%%%%%%%%%%%%%%%%%
\begin{itemize}
\item
	\ressubheading{Spectral Heuristics for $k$-Incoherent Decompositions}{}{Atlantic Undergraduate Physics and Astronomy Conference 2026 (UPEI)}{March 2026}

\item
	\ressubheading{Advancements in Graph Bandwidth Reduction}{}{Science Atlantic MSCS Conference 2025 (Cape Breton University)}{October 2025}

\item
	\ressubheading{$S$-Bandwidth as an Indicator of PST on Quantum Networks}{}{Atlantic Undergraduate Physics and Astronomy Conference 2025 (MUN)}{February 2025}

\item
	\ressubheading{Computing the $S$-Bandwidth of a Quantum Network}{}{Science Atlantic MSCS Conference 2024 (Acadia University)}{October 2024}
\end{itemize}


%%%%%%%%%%%%%%%%%%%%%%%%%%%%%%
\resheading{Research Experience}
%%%%%%%%%%%%%%%%%%%%%%%%%%%%%%
\begin{itemize}
\item
	\ressubheading{Student Researcher}{Mount Allison University}{Computer Science}{May 2026 -- Aug. 2026}
	\begin{itemize}
		\resitem{Investigated the training dynamics of bandwidth-constrained RNNs using Python/PyTorch}
        \begin{itemize}
            \resitem{Ran formal language classification and modelling experiments using DFA-sampled strings}
            \resitem{Tracked and analyzed eigenvalue geometry to explain observed differences in trainability}
            \resitem{Wrote custom CUDA kernels and CPU routines in C for inference time benchmarking}
        \end{itemize}
        \resitem{\textbf{Advised by:} Dr.~Nathaniel Johnston, MtA Department of Mathematics \& Computer Science}
	\end{itemize}

\item
	\ressubheading{Student Researcher}{Mount Allison University}{Political Science}{Apr. 2026 -- Aug. 2026}
	\begin{itemize}
		\resitem{Investigated gaps in political philosophy frameworks given modern technological developments}
        \resitem{\textbf{Advised by:} Dr.~Wayne A.~Hunt, MtA Department of Politics \& International Relations}
	\end{itemize}

\item
	\ressubheading{Research Assistant}{Mount Allison University}{Economics}{Sep. 2025 -- Dec. 2025}
	\begin{itemize}
		\resitem{Modelled the effects of exogenous expenditure on short-term municipal tax rates using Python}
        \resitem{\textbf{Worked for:} Dr.~Craig Brett, MtA Department of Economics}
	\end{itemize}

\item
	\ressubheading{Student Researcher}{Government of New Brunswick}{Public Policy}{May 2025 -- Aug. 2025}
	\begin{itemize}
		\resitem{Studied limitations of accessibility policy/legislation in the New Brunswick education system}
        \resitem{\textbf{Worked for:} Accessibility Office, NB Department of Post-Secondary, Training and Labour}
	\end{itemize}

\item
	\ressubheading{Research Assistant}{Mount Allison University}{Mathematics}{Apr. 2024 -- Aug. 2024; May 2025 -- Aug. 2025}
	\begin{itemize}
        \resitem{Explored tightened Cauchy--Schwarz-type bounds and generalized to matrices/multiple vectors}
        \resitem{Ran computational surveys of Laplacian integral and $\{-1,0,1\}$-diagonalizable graphs in Julia}
        \resitem{Implemented various algorithms in MATLAB to assist with quantum information research}
        \resitem{\textbf{Worked for:} Dr.~Nathaniel Johnston, MtA Department of Mathematics \& Computer Science}
	\end{itemize}
\end{itemize}


%%%%%%%%%%%%%%%%%%%%%%%%%%%%%%
\resheading{Teaching Experience}
%%%%%%%%%%%%%%%%%%%%%%%%%%%%%%
\begin{itemize}
\item 
	\ressubheading{Calculus II (MATH 1121)}{Mount Allison University}{Teaching Assistant}{Winter 2026}
	\begin{itemize}
		\resitem{Co-led weekly labs with $30+$ students, fielding questions and marking half of the assignments}
	\end{itemize}

\item 
	\ressubheading{Math Help Centre}{Mount Allison University}{Teaching Assistant}{Winter 2025, Winter 2026}
	\begin{itemize}
		\resitem{Supervised semiweekly sessions with students from $6+$ different first- and second-year courses}
	\end{itemize}

\item 
	\ressubheading{Linear Algebra (MATH 2221)}{Mount Allison University}{Marker}{Winter 2025}
	\begin{itemize}
		\resitem{Marked $30+$ student assignments per week}
	\end{itemize}

\item 
	\ressubheading{Applied Calculus (MATH 1151)}{Mount Allison University}{Teaching Assistant}{Fall 2024}
	\begin{itemize}
		\resitem{Co-led weekly labs with $30+$ students, fielding questions and marking half of the assignments}
	\end{itemize}
\end{itemize}


%%%%%%%%%%%%%%%%%%%%%%%%%%%%%%
\resheading{Awards, Grants, and Honours}
%%%%%%%%%%%%%%%%%%%%%%%%%%%%%%
	\vspace{-10pt}
	\begin{center}\begin{longtable}{l@{\extracolsep{\fill}}r}
		\multicolumn{2}{c}{Independent Student Research Grant (\$$9,\!000$) --- Mount Allison University \cftdotfill{\cftdotsep} 2026}\\
		\multicolumn{2}{c}{$3$\textsuperscript{rd} Place Team Award --- Atlantic Canada Programming Competition \cftdotfill{\cftdotsep}2025}\\
		\multicolumn{2}{c}{$2$\textsuperscript{nd} Place Undergraduate Research Award for Computer Science --- Science Atlantic \cftdotfill{\cftdotsep}2025}\\
		\multicolumn{2}{c}{President's Scholarship --- Mount Allison University \cftdotfill{\cftdotsep} 2022} \\
		\vphantom{E}
\end{longtable}
\end{center}\vspace*{-30pt}


%%%%%%%%%%%%%%%%%%%%%%%%%%%%%%
\resheading{Other Academic and Professional Activities}
%%%%%%%%%%%%%%%%%%%%%%%%%%%%%%
\begin{itemize}
\item
	Creator and maintainer of \hyperref{https://luis-varona.github.io/MatrixBandwidth.jl/}{}{}{MatrixBandwidth.jl}, a peer-reviewed Julia package offering fast algorithms for matrix bandwidth minimization and recognition. \\
{\em June 2025 -- present}

\item
	Contributor to the \hyperref{http://www.oeis.org}{}{}{On-Line Encyclopedia of Integer Sequences}. \\
	{\em April 2025 -- present}

\item
	Core member of the Mount Allison Competitive Programming Team. \\
{\em November 2024 -- present}

\item
	Contributed to the following open-source software packages:{}\vspace{-0.05in}
	\begin{itemize}
		\item \hyperref{http://www.qetlab.com}{}{}{QETLAB}, a MATLAB package for exploring quantum entanglement.
        \item \hyperref{https://github.com/spectregrams/spectre/wiki}{}{}{Spectre}, an SDR-agnostic Python program for recording radio signals and spectrograms.
        \item \hyperref{https://scikit-learn.org}{}{}{scikit-learn}, a machine learning library for Python.
        \item \hyperref{https://pola.rs/}{}{}{Polars}, a high-performance DataFrame library for Python and Rust
        \item The \hyperref{https://julialang.org/}{}{}{Julia} standard library (in particular, \hyperref{https://github.com/JuliaLang/LinearAlgebra.jl}{}{}{LinearAlgebra.jl}).
	\end{itemize}
\end{itemize}
\end{document}
